\documentclass[../main.tex]{subfiles}
\begin{document}
\section{Scope of an encrypted-search application}
\label{sec:trapdoor_rel_model}
Approximate relations can model the \emph{functionality} of a private search
index: tuples may encode document--term incidence, co-occurrence, or candidate
matches, and a positive query combines approximate probes. The source map must
still reflect repeated trapdoors, common stored filters, and canonical keys.
Subject to a justified construction law, the preceding bounds describe false
matches and missed answers. They can guide allocation of an error budget
before a separate security analysis.

This functionality model does not specify confidentiality. Two structures
with identical FP and FN rates can expose different sizes, access patterns,
query-equality patterns, or stored representations. A secret salt alone does
not establish an encrypted-search security theorem. Such a theorem needs a
concrete construction, adversary, experiment, and declared leakage, including
what repeated and adaptive queries reveal. In particular, the probability
of a false positive on an exact nonmember is not the posterior probability
that a returned match is false. No cryptographic security claim follows from
the Bernoulli channel parameters or from the experiments in this paper.
\end{document}
